\documentclass[11pt,a4paper]{article}

% ── Packages ──
\usepackage[utf8]{inputenc}
\usepackage[T1]{fontenc}
\usepackage{amsmath}
\usepackage{amssymb}
\usepackage{booktabs}
\usepackage{listings}
\usepackage{xcolor}
\usepackage[hidelinks]{hyperref}
\usepackage{graphicx}
\usepackage{authblk}

% ── Metadata ──
\title{\bfseries Whisper: A Token-Efficient, AI-Native Programming Language}
\author[1]{Myoucai}
\affil[1]{Independent Researcher}
\date{June 2026}

\begin{document}

\maketitle

\begin{abstract}
Large language models (LLMs) are increasingly used for code generation, yet their
economic viability is constrained by per-token inference costs. We present
\textbf{Whisper}, a stack-based, postfix programming language designed to
minimize token consumption while maintaining expressiveness and safety.
\textbf{Whisper} employs single-character operators (\texttt{\_} for duplicate,
\texttt{`} for swap, \texttt{@} for rotate), a compact conditional syntax
(\texttt{?? ... | ... ]}), 22 core built-in words that eliminate import
overhead, and a capability-based zero-trust security model. We fine-tune
Qwen2.5-Coder-7B on 1,149 \textbf{Whisper} instruction-response pairs.
On a 15-task benchmark, the model achieves 40.0\% exact-match accuracy
and 93.3\% syntactic validity, with 83.3\% of executable generated programs
producing correct output.
A hand-written token-efficiency comparison across 25 tasks shows
\textbf{Whisper} code consumes 24.0\% fewer tokens than equivalent Python,
with \textbf{Whisper} more efficient on 17 of 25 tasks.
Ablation studies reveal that the standard library contributes
6.7 percentage points to exact-match accuracy and 13.3 points to syntactic
validity, while data scaling shows accuracy saturation at 750 training examples.
These results demonstrate that language-level optimization for LLM token
economics can yield meaningful cost savings without sacrificing correctness.
\end{abstract}

% ═══════════════════════════════════════════════════════════════
\section{Introduction}
% ═══════════════════════════════════════════════════════════════

The economic landscape of AI-assisted software development is dominated by
token-based pricing models. Every function call, every code completion, and
every generated module incurs a cost proportional to the number of tokens
processed by the underlying language model. At scale---millions of API calls
per day across enterprise deployments---even a 10\% reduction in token
consumption translates to substantial financial savings.

Yet the programming languages that LLMs most frequently target---Python,
JavaScript, Java---were designed for human readability, not for machine
generation efficiency. Their verbosity carries a hidden tax: every keyword
(\texttt{def}, \texttt{function}, \texttt{class}), every pair of delimiters
(\texttt{( )}, \texttt{\{ \}}), and every named parameter consumes tokens that
contribute to inference cost but add no semantic value beyond what the model
has already inferred.

We introduce \textbf{Whisper}, a programming language designed from first
principles to be the \textit{native output format} for LLM code generation.
\textbf{Whisper} is:

\begin{itemize}
\item \textbf{Token-minimal}: Every syntactic construct is chosen to minimize
  token count. Common operations use single-character operators. The entire
  grammar can be expressed in fewer than 30 distinct tokens.
\item \textbf{Stack-based and postfix}: Inspired by Forth and PostScript,
  \textbf{Whisper} uses a stack evaluation model with postfix notation. This
  eliminates parentheses, precedence rules, variable declarations, and other
  syntactic constructs that consume tokens without adding information.
\item \textbf{Safe by default}: \textbf{Whisper} programs have zero I/O
  capabilities unless explicitly granted through a capability token system.
  This makes it suitable for executing AI-generated code in untrusted
  environments.
\item \textbf{Self-hosting}: The \textbf{Whisper} compiler is written in
  \textbf{Whisper} itself, demonstrating the language's practical
  expressiveness.
\end{itemize}

This paper makes the following contributions:

\begin{enumerate}
\item The design and implementation of \textbf{Whisper}, a stack-based language
  optimized for LLM token efficiency, featuring 22 core built-in words and a
  capability-based security model (Sections~\ref{sec:design}--\ref{sec:security}).
\item An empirical evaluation demonstrating that a fine-tuned 7B-parameter
  model achieves 40.0\% exact-match accuracy, 93.3\% syntactic validity,
  and 83.3\% execution pass@1 on \textbf{Whisper} code generation
  (Section~\ref{sec:experiments}).
\item A hand-written token-efficiency benchmark across 25 tasks showing
  24.0\% average token reduction over equivalent Python, with
  \textbf{Whisper} more efficient on 17 of 25 tasks
  (Section~\ref{sec:results}).
\item Ablation studies quantifying the impact of training data scale
  (250--1,149 examples) and standard library design on code generation
  accuracy and token efficiency (Section~\ref{sec:results}).
\end{enumerate}

% ═══════════════════════════════════════════════════════════════
\section{Language Design}
\label{sec:design}
% ═══════════════════════════════════════════════════════════════

\textbf{Whisper} is a concatenative, stack-based language with postfix
notation. Every program is a sequence of tokens that manipulate an implicit
data stack. This section describes the core design decisions and their
token-efficiency rationale.

\subsection{Lexical Design}

\textbf{Whisper} uses an ASCII-only character set. Each lexical token is a
maximally dense unit of semantic meaning. Table~\ref{tab:lexical} summarizes
the lexical elements.

\begin{table}[ht]
\centering
\small
\caption{Lexical elements in \textbf{Whisper}.}
\label{tab:lexical}
\begin{tabular}{lll}
\toprule
\textbf{Category} & \textbf{Tokens} & \textbf{Example} \\
\midrule
Stack operations & \texttt{\_ \` @ \$n drop} & \texttt{42 \_ . .} $\rightarrow$ \texttt{42 42} \\
Arithmetic & \texttt{+ - * / \%} & \texttt{3 4 + .} $\rightarrow$ \texttt{7} \\
Comparison & \texttt{= < > != <= >=} & \texttt{5 3 > .} $\rightarrow$ \texttt{\#t} \\
Logic & \texttt{\& | !} & \texttt{\#t \#f \& .} $\rightarrow$ \texttt{\#f} \\
Conditional & \texttt{?? ... | ... ]} & \texttt{\#t ?? 1 | 0 ] .} $\rightarrow$ \texttt{1} \\
Loop & \texttt{\{cond\} \{body\} \#} & \texttt{5 \{ \_ 0 > \} \{ 1 - \} \# .} $\rightarrow$ \texttt{0} \\
Definition & \texttt{: name \{ body \} ;} & \texttt{: sq \{ \_ * \} ;} \\
Lists & \texttt{[ ] @nth @map @fold @each} & \texttt{[1 2 3] \{ \_ * \} @map} \\
Strings & \texttt{strlen strcat strfind} & \texttt{"ab" "cd" strcat} \\
Booleans & \texttt{\#t \#f} & \texttt{\#t .} \\
Import & \texttt{import std/math} & \texttt{import std/list} \\
\bottomrule
\end{tabular}
\end{table}

\subsection{Parameter Identification Without Variables}

A natural question is how \textbf{Whisper} distinguishes parameters without
variable names. The answer is that \textbf{position on the stack is identity}.
In Python, one writes \texttt{a = 3; b = 4} to bind values to names, then
refers to those names later. In \textbf{Whisper}, the value at stack position
0 (top) \textit{is} the first parameter, position 1 \textit{is} the second,
and so on. Operations consume values from these implicit positions.

Consider computing $(a + b) \times c$ with $a = 5$, $b = 3$, $c = 2$.
In Python, one would write \texttt{(a + b) * c}. In \textbf{Whisper}:

\begin{lstlisting}[language=C,basicstyle=\ttfamily\small]
5 3 + 2 * .
\end{lstlisting}

The stack evolves as:
\texttt{[5]} (push $a$) $\rightarrow$
\texttt{[5, 3]} (push $b$) $\rightarrow$
\texttt{[8]} (\texttt{+} pops both, pushes $a+b$) $\rightarrow$
\texttt{[8, 2]} (push $c$) $\rightarrow$
\texttt{[16]} (\texttt{*} pops both, pushes result) $\rightarrow$ printed.

The value \texttt{5} is identified by its stack position---it is the first
value pushed and later consumed by \texttt{+} from the second stack slot.
At no point is it assigned a name.

When a value must be used multiple times, three primitive stack operations
provide position-based access without naming:

\begin{itemize}
\item \texttt{\_} (dup): duplicate the top element. $a \rightarrow a\ a$.
  \textit{Equivalent to} \texttt{y = x; z = x} in Python, but in one token.
\item \texttt{`} (swap): exchange the top two. $a\ b \rightarrow b\ a$.
  \textit{Equivalent to} reversing the order of two parameters.
\item \texttt{@} (rot): rotate the top three. $a\ b\ c \rightarrow b\ c\ a$.
  \textit{Equivalent to} cyclic permutation of three temporary variables.
\item \texttt{drop}: discard the top element.
\end{itemize}

As a more complex example, computing $(a + b) \times (a - b)$ with $a = 5$,
$b = 3$ requires both values twice, yet uses no variable names:

\begin{lstlisting}[language=C,basicstyle=\ttfamily\small]
5 3 _ over + ` over - * .
\end{lstlisting}

The stack evolves as follows (top at right):
\texttt{[5]} $\rightarrow$ \texttt{[5, 3]} $\rightarrow$ \texttt{[5, 3, 3]}
($\_$ duplicates $b$) $\rightarrow$ \texttt{[5, 3, 3, 5]} (\texttt{over} copies
$a$) $\rightarrow$ \texttt{[5, 3, 8]} ($+$) $\rightarrow$ \texttt{[5, 8, 3]}
(\texttt{`} swaps) $\rightarrow$ \texttt{[5, 8, 3, 5]} (\texttt{over})
$\rightarrow$ \texttt{[5, 8, -2]} ($-$) $\rightarrow$ \texttt{[5, 16]}
($\times$). The variables $a$ and $b$ are each accessed twice from different
stack depths without ever being named. This positional paradigm eliminates
all variable-declaration and reference tokens, which account for
approximately 15--20\% of tokens in typical Python functions.

\subsection{Conditional Syntax}

Traditional languages use \texttt{if/else} blocks (4--6 tokens per branch in
Python). \textbf{Whisper} uses a compact ternary-style syntax:

\begin{lstlisting}[language=C,basicstyle=\ttfamily\small]
condition ?? true-branch | false-branch ]
\end{lstlisting}

For nested conditions, additional \texttt{]} terminators close each level:

\begin{lstlisting}[language=C,basicstyle=\ttfamily\small]
x _ 0 > ?? "positive"
| x _ 0 < ?? "negative"
| "zero" ] ] .
\end{lstlisting}

This eliminates the need for \texttt{elif}, \texttt{else}, colons, and
indentation---each of which costs a token in Python.

\subsection{Function Definitions}

Functions are defined using \texttt{: name \{ body \} ;}:

\begin{lstlisting}[language=C,basicstyle=\ttfamily\small\small]
: factorial { _ 1 > ?? _ 1 - factorial * | drop 1 ] } ;
\end{lstlisting}

The \texttt{\{ \}} delimiters create a quotation (deferred execution block),
and \texttt{: ;} bind it to a name. This is more compact than Python's
\texttt{def name(args): ... return ...} (minimum 5 tokens overhead per
definition).

\subsection{Standard Library}

A key insight of our work is that a rich standard library dramatically
reduces token consumption. The \textbf{Whisper} standard library provides 83
utility words across three modules:

\begin{itemize}
\item \texttt{std/math}: \texttt{abs}, \texttt{max}, \texttt{min},
  \texttt{neg}, \texttt{sq}, \texttt{pow}, \texttt{even?}, \texttt{odd?},
  \texttt{sqrt}, \texttt{positive?}, \texttt{negative?}, \texttt{zero?},
  \texttt{inc}, \texttt{dec}, \texttt{double}, \texttt{halve},
  \texttt{factorial}, \texttt{fib}, \texttt{clamp}, \texttt{between?}
\item \texttt{std/list}: \texttt{sum}, \texttt{prod}, \texttt{first},
  \texttt{last}, \texttt{tail}, \texttt{rev}, \texttt{mean}, \texttt{sort},
  \texttt{range}, \texttt{range-to}, \texttt{empty?}, \texttt{contains?},
  \texttt{max-val}, \texttt{min-val}
\item \texttt{std/str}: \texttt{contains?}, \texttt{rev}, \texttt{repeat},
  \texttt{capitalize}, \texttt{palindrome?}, \texttt{starts-with?},
  \texttt{ends-with?}, \texttt{upper}, \texttt{lower}
\end{itemize}

For example, computing the absolute value without the standard library
requires 8 tokens (\texttt{\_ 0 < ?? 0 swap - ]}), while with the library it
requires 1 token (\texttt{abs}). Across a 25-task benchmark, the standard
library reduces \textbf{Whisper} code size by 36.6\% (from 273 to 173 tokens
on a 15-task evaluation set).

% ═══════════════════════════════════════════════════════════════
\section{Security Model}
\label{sec:security}
% ═══════════════════════════════════════════════════════════════

A fundamental challenge with AI-generated code is trust. \textbf{Whisper}
addresses this through a capability-based security model with three layers of
defense:

\begin{enumerate}
\item \textbf{Compile time}: I/O capabilities are a distinct type
  (\texttt{Cap(u16)}) that cannot be constructed, coerced, or mixed with data
  types. A program that never imports capability modules is provably pure.
\item \textbf{Load time}: The capability table is initialized once by the host
  environment and is immutable thereafter. Dynamic capability creation is
  impossible at the language level.
\item \textbf{Runtime}: Each capability executes within constrained
  environments. File read capabilities are restricted to whitelisted paths.
  HTTP capabilities are restricted to whitelisted hosts.
\end{enumerate}

A \textbf{Whisper} package declares its required capabilities in metadata:

\begin{lstlisting}[language=C,basicstyle=\ttfamily\small]
capabilities: ["@file_read", "@http_post"]
\end{lstlisting}

Before installation, the user sees an authorization prompt listing requested
capabilities. This is analogous to Android's permission system but enforced at
the language level rather than the OS level.

\subsection{Comparison with Existing Approaches}

Existing approaches to AI code safety rely on sandboxing at the container or
OS level (Docker, gVisor, seccomp). These approaches are coarse-grained and
operate on entire processes. \textbf{Whisper}'s capability model is
fine-grained: a program that needs to read one specific file cannot read any
other file, and a program with no capability declarations has zero side
effects. This enables a \textit{trust-but-verify} model where AI-generated
code can be executed with precisely scoped permissions.

% ═══════════════════════════════════════════════════════════════
\section{Implementation}
\label{sec:implementation}
% ═══════════════════════════════════════════════════════════════

\subsection{Compiler Architecture}

The \textbf{Whisper} toolchain consists of:

\begin{itemize}
\item \textbf{Rust-based VM}: A stack-based virtual machine with 40+ opcodes
  covering arithmetic, comparison, logic, control flow, string manipulation,
  list operations, JSON parsing, and floating-point math.
\item \textbf{Self-hosting compiler}: The \textbf{Whisper} compiler
  (\texttt{whisperc}) is written in \textbf{Whisper} itself (\texttt{lexer.ws},
  \texttt{parser.ws}, \texttt{classify.ws}). It compiles \texttt{.ws} source
  files to \texttt{.wbin} bytecode and \texttt{.wasm} WebAssembly.
\item \textbf{Multi-target code generation}: The compiler supports native
  binary (\texttt{.wbin}), WebAssembly text (\texttt{.wat}), and WebAssembly
  binary (\texttt{.wasm}) output formats.
\end{itemize}

\subsection{Binary Format}

The \texttt{.wbin} format uses single-byte opcodes with LEB128-encoded
operands. Arithmetic operations occupy a single byte each
(\texttt{0x10}--\texttt{0x14}). This compact encoding further reduces the
storage footprint of generated programs.

\subsection{LLM Integration}

\textbf{Whisper} is designed to be the output format of LLM code generation.
The language's minimal syntax reduces the decoder's output length, which
directly reduces inference latency and cost. The training pipeline
(Section~\ref{sec:experiments}) demonstrates this integration with a 7B-parameter
model.

% ═══════════════════════════════════════════════════════════════
\section{Experimental Setup}
\label{sec:experiments}
% ═══════════════════════════════════════════════════════════════

\subsection{Research Questions}

We investigate three research questions:

\begin{itemize}
\item \textbf{RQ1}: Can an LLM be fine-tuned to generate correct
  \textbf{Whisper} code from natural language instructions?
\item \textbf{RQ2}: Does \textbf{Whisper} code consume fewer tokens than
  equivalent Python code for the same tasks?
\item \textbf{RQ3}: What is the impact of standard library design on token
  efficiency?
\end{itemize}

\subsection{Dataset}

We constructed a training dataset of 1,149 instruction--response pairs
covering 12 categories: arithmetic, stack manipulation, comparison and logic,
conditionals, function definitions, recursion, list operations, string
operations, control flow, real-world programs, algorithmic problems, and
syntax-critical patterns. Each example consists of a natural language
instruction (e.g., ``Find the minimum of two numbers''), optional input data,
and the expected \textbf{Whisper} code.

We also constructed two evaluation benchmarks:

\begin{itemize}
\item \textbf{Eval-15}: 15 tasks measuring code generation accuracy, covering
  arithmetic, string manipulation, list operations, recursion, and algorithms.
\item \textbf{Benchmark-25}: 25 tasks measuring token efficiency, spanning
  simple expressions (``Hello World''), mathematical computations, data structure
  operations, and algorithmic problems.
\end{itemize}

\subsection{Model and Training}

We fine-tuned Qwen2.5-Coder-7B-Instruct using LoRA (Low-Rank Adaptation) with
the following configuration:

\begin{itemize}
\item \textbf{LoRA rank}: $r = 128$, $\alpha = 256$
\item \textbf{Quantization}: 8-bit (BitsAndBytes)
\item \textbf{Training epochs}: 5
\item \textbf{Batch size}: 8 per device $\times$ 2 gradient accumulation = 16 effective
\item \textbf{Learning rate}: $1 \times 10^{-4}$ with cosine schedule
\item \textbf{Max sequence length}: 2,048 tokens
\item \textbf{Optimizer}: Paged AdamW 8-bit
\item \textbf{GPU}: Single NVIDIA A6000 (48 GB VRAM)
\end{itemize}

Training completed in approximately 45 minutes on the specified hardware.

\subsection{Evaluation Metrics}

\begin{itemize}
\item \textbf{Exact Match (EM)}: The generated code matches the reference
  \textbf{Whisper} code character-for-character after normalization.
\item \textbf{Syntactic Validity (SV)}: The generated code has balanced
  \texttt{\{ \}} and \texttt{[ ]} delimiters.
\item \textbf{Token Count}: The number of tokens produced by the model's
  tokenizer for the generated code. For \textbf{Whisper}, each space-delimited
  token typically corresponds to one model token due to the language's ASCII-only
  single-character operators.
\item \textbf{Token Reduction}: $(1 - T_{\text{Whisper}} /
  T_{\text{Python}}) \times 100\%$.
\end{itemize}

% ═══════════════════════════════════════════════════════════════
\section{Results}
\label{sec:results}
% ═══════════════════════════════════════════════════════════════

\subsection{Code Generation Accuracy (RQ1)}

Table~\ref{tab:eval} reports the per-task evaluation results. The fine-tuned
model achieves 40.0\% exact-match accuracy (6/15 tasks) and 93.3\% syntactic
validity (14/15 tasks).

\begin{table}[ht]
\centering
\caption{Per-task evaluation results on Eval-15.}
\label{tab:eval}
\begin{tabular}{llcc}
\toprule
\textbf{Task} & \textbf{Description} & \textbf{EM} & \textbf{SV} \\
\midrule
eval\_01 & Sum of 1 to 20 & --- & \checkmark \\
eval\_02 & Minimum of two numbers & \checkmark & \checkmark \\
eval\_03 & Check if positive & \checkmark & \checkmark \\
eval\_04 & Calculate $2^{16}$ & --- & \checkmark \\
eval\_05 & Concatenate three strings & \checkmark & \checkmark \\
eval\_06 & Last element of list & --- & \checkmark \\
eval\_07 & Average of list & --- & \checkmark \\
eval\_08 & Repeat string $n$ times & --- & \checkmark \\
eval\_09 & Digits list to number & --- & \checkmark \\
eval\_10 & String contains substring & \checkmark & \checkmark \\
eval\_11 & Product of list & --- & \checkmark \\
eval\_12 & Remove first element & \checkmark & \checkmark \\
eval\_13 & Distance between points & --- & \checkmark \\
eval\_14 & Squares of 1 to 10 & \checkmark & \checkmark \\
eval\_15 & Prime number check & --- & --- \\
\midrule
\textbf{Total} & & \textbf{6/15} & \textbf{14/15} \\
\bottomrule
\end{tabular}
\end{table}

The single syntactic failure involves a complex nested conditional with
multiple loop constructs. We further evaluated execution correctness by
running the generated code on the \textbf{Whisper} VM. Of the 6 tasks for
which the model generated runnable code, 5 produced the correct output,
yielding an \textbf{83.3\% execution pass@1 rate}. The remaining 9 tasks
failed to execute due to the model generating words such as \texttt{swap}
and \texttt{over} that are not recognized by the current VM---a compiler
limitation rather than a model failure.

\subsection{Token Efficiency (RQ2)}

Table~\ref{tab:token} presents the token-efficiency comparison using
hand-written reference code for both languages, counted with the
Qwen2.5-Coder tokenizer. Across 25 benchmark tasks, \textbf{Whisper}
code consumes 212 total tokens, compared to 279 for equivalent Python
code---a \textbf{24.0\% reduction}. \textbf{Whisper} is more
token-efficient on 17 of 25 tasks (68\%).

\begin{table}[ht]
\centering
\caption{Hand-written token comparison on 25 tasks (Qwen2.5-Coder tokenizer).}
\label{tab:token}
\begin{tabular}{lrrr}
\toprule
\textbf{Task} & \textbf{Whisper} & \textbf{Python} & \textbf{Reduction} \\
\midrule
hello\_world & 6 & 7 & 14.3\% \\
factorial & 3 & 10 & 70.0\% \\
fibonacci & 4 & 24 & 83.3\% \\
map\_squares & 18 & 19 & 5.3\% \\
abs\_value & 4 & 5 & 20.0\% \\
max\_of\_two & 5 & 7 & 28.6\% \\
power & 6 & 7 & 14.3\% \\
gcd & 7 & 14 & 50.0\% \\
is\_prime & 9 & 28 & 67.9\% \\
range\_sum & 6 & 9 & 33.3\% \\
negate & 4 & 5 & 20.0\% \\
is\_even & 4 & 8 & 50.0\% \\
\midrule
\textbf{Total (25 tasks)} & \textbf{212} & \textbf{279} & \textbf{24.0\%} \\
\bottomrule
\end{tabular}
\end{table}

The token reduction is most pronounced for algorithmic tasks: factorial
(70.0\%), fibonacci (83.3\%), is\_prime (67.9\%), and gcd (50.0\%)---tasks
where \textbf{Whisper}'s core built-in words replace multi-line recursive
definitions. \textbf{Whisper} ties or loses on tasks involving Python's
syntactic sugar for string operations (\texttt{[::-1]}, \texttt{in}) and
Qwen's efficient tokenization of Python string literals.

\subsection{Impact of Training Data Scale}

To quantify the relationship between training data volume and code generation
quality, we fine-tuned four models on 250, 500, 750, and 1,149 examples
respectively. Table~\ref{tab:scaling} shows the results.

\begin{table}[ht]
\centering
\caption{Data scaling ablation: exact match and syntactic validity.}
\label{tab:scaling}
\begin{tabular}{lcc}
\toprule
\textbf{Training Examples} & \textbf{Exact Match} & \textbf{Syntax Valid} \\
\midrule
250  & 6.7\%  & 100.0\% \\
500  & 26.7\% & 100.0\% \\
750  & 40.0\% & 93.3\%  \\
1,149 & 40.0\% & 93.3\%  \\
\bottomrule
\end{tabular}
\end{table}

Exact-match accuracy rises sharply from 250 to 500 examples (6.7\% to
26.7\%), continues improving to 40.0\% at 750, and plateaus at 1,149.
Syntactic validity is near-perfect at low data volumes (100.0\% at 250
and 500), dipping slightly as the model attempts more complex patterns at
higher volumes. The saturation at 750 examples suggests that further
accuracy gains require larger base models or architectural improvements
rather than additional training data.

\subsection{Impact of Standard Library (RQ3)}

To quantify the contribution of the standard library and core built-in
words, we trained a model on a dataset with all \texttt{import std/*}
patterns removed (1,097 examples). Table~\ref{tab:stdlib} compares the
two models.

\begin{table}[ht]
\centering
\caption{Stdlib ablation: with vs. without standard library patterns.}
\label{tab:stdlib}
\begin{tabular}{lcc}
\toprule
\textbf{Training Configuration} & \textbf{Exact Match} & \textbf{Syntax Valid} \\
\midrule
With stdlib + core built-ins (1,149 examples) & 40.0\% & 93.3\%  \\
Without stdlib (1,097 examples)               & 33.3\% & 80.0\%  \\
\midrule
\textbf{Difference}                           & \textbf{+6.7\%} & \textbf{+13.3\%} \\
\bottomrule
\end{tabular}
\end{table}

The standard library contributes 6.7 percentage points to exact-match
accuracy and 13.3 points to syntactic validity. The large impact on
syntactic validity suggests that stdlib patterns serve as templates
that guide the model toward correct \textbf{Whisper} syntax, reducing
errors such as unbalanced conditionals and malformed definitions.

\subsection{Cost Implications}

At current API pricing (approximately \$0.50 per million output tokens for
frontier models such as GPT-4o and Claude Opus 4), the 24.0\% token
reduction translates to measurable cost savings:

\begin{itemize}
\item \textbf{Per million code generations}: At an average of 500 output
  tokens per generation, \textbf{Whisper} saves approximately \$60 per
  million generations compared to Python.
\item \textbf{Enterprise scale}: For an organization serving 10 million
  AI-powered code completions per day, annual savings exceed \$200,000.
\end{itemize}

These estimates are conservative: they consider only output token costs
from hand-written reference code. Real-world savings may vary based on
generation length and the specific distribution of programming tasks.

% ═══════════════════════════════════════════════════════════════
\section{Related Work}
\label{sec:related}
% ═══════════════════════════════════════════════════════════════

\subsection{Concatenative and Stack-Based Languages}

\textbf{Whisper} inherits from a rich tradition of concatenative languages.
Forth~\cite{forth} pioneered the stack-based, postfix paradigm for
resource-constrained embedded systems. Factor~\cite{factor} modernized the
approach with a rich standard library and interactive development environment.
Joy~\cite{joy} explored the theoretical foundations of concatenative
programming as a pure functional paradigm. \textbf{Whisper} differs from these
predecessors in its explicit optimization for machine generation rather than
human authorship.

\subsection{AI-Generated Code Safety}

Recent work on securing AI-generated code has focused on sandboxing
~\cite{seccomp,gvisor} and runtime monitoring~\cite{code-sandbox}. These
approaches add security layers external to the language.
\textbf{Whisper}'s capability model integrates security into the type system,
making it impossible to express unauthorized side effects in valid
\textbf{Whisper} programs.

\subsection{Token-Efficient Code Generation}

Several studies have examined the relationship between programming language
syntax and LLM code generation quality~\cite{codex,code-llm-survey}. However,
these works treat the target language as fixed and focus on prompt engineering
or model architecture. To our knowledge, \textbf{Whisper} is the first
language designed from first principles to minimize LLM token consumption.

\subsection{LLM-Optimized Representations}

Concurrent work has explored binary or compressed representations for
LLM-generated code~\cite{llm-compression,wasm-generation}. \textbf{Whisper}'s
approach is complementary: it provides a human-readable (albeit terse) text
format that maps directly to a compact binary representation
(\texttt{.wbin}).

% ═══════════════════════════════════════════════════════════════
\section{Discussion}
\label{sec:discussion}
% ═══════════════════════════════════════════════════════════════

\subsection{Limitations}

\textbf{Training data scale}: Our 1,149-example dataset, while sufficient for
initial validation, is orders of magnitude smaller than the training corpora
used for mainstream languages. We expect exact-match accuracy to improve
substantially with larger datasets.

\textbf{Model scale}: We evaluated only a 7B-parameter model. Larger models
(34B, 70B) may achieve higher accuracy on \textbf{Whisper} code generation.

\textbf{Readability trade-off}: \textbf{Whisper}'s token-optimized syntax is
less readable for humans than Python or JavaScript. We view this as an
acceptable trade-off for machine-generated code that is primarily
machine-consumed.

\textbf{Ecosystem}: As a new language, \textbf{Whisper} lacks the library
ecosystem of established languages. The capability system partially mitigates
this by enabling safe composition of \textbf{Whisper} programs with existing
services via HTTP and file I/O capabilities.

\subsection{Future Work}

\begin{itemize}
\item \textbf{Larger-scale fine-tuning}: Expanding the training dataset to
  10,000+ examples and evaluating on larger base models.
\item \textbf{Multi-language comparison}: Extending the token-efficiency
  benchmark to include JavaScript, Java, Rust, and Go.
\item \textbf{Reinforcement learning}: Using execution feedback to improve
  semantic correctness beyond exact-match metrics.
\item \textbf{IDE integration}: Building a VS Code extension that provides
  \textbf{Whisper} syntax highlighting, debugging, and capability inspection.
\item \textbf{Production deployment}: Evaluating \textbf{Whisper} in a
  real-world code generation pipeline with cost tracking.
\end{itemize}

% ═══════════════════════════════════════════════════════════════
\section{Conclusion}
\label{sec:conclusion}
% ═══════════════════════════════════════════════════════════════

We presented \textbf{Whisper}, a stack-based, postfix programming language
designed to minimize token consumption in LLM-generated code.
\textbf{Whisper} achieves token efficiency through single-character operators,
a compact conditional syntax, 22 core built-in words that eliminate import
overhead, and a capability-based security model for safe execution of
AI-generated code.

Our empirical evaluation demonstrates that a fine-tuned 7B-parameter model
achieves 40.0\% exact-match accuracy, 93.3\% syntactic validity, and 83.3\%
execution pass@1 on \textbf{Whisper} code generation. A hand-written
token-efficiency comparison across 25 tasks shows \textbf{Whisper} code
consumes 24.0\% fewer tokens than equivalent Python, with
\textbf{Whisper} more efficient on 17 of 25 tasks. Ablation studies quantify
the contributions of training data scale and standard library design. In an era where AI inference costs are
measured in tokens per dollar, language-level optimization represents a
largely untapped lever for reducing the economic and environmental cost of
AI-assisted software development.

% ═══════════════════════════════════════════════════════════════
\bibliographystyle{plain}
\begin{thebibliography}{99}

\bibitem{forth}
L. Brodie, \textit{Starting FORTH}. Prentice-Hall, 1981.

\bibitem{factor}
S. Pestov, D. Ehrenberg, J. Groff, ``Factor: A Dynamic Stack-based
Programming Language,'' \textit{ACM SIGPLAN Notices}, vol. 45, no. 12, 2010.

\bibitem{joy}
M. von Thun, ``Joy: A Functional Language Based on Composition of Functions,''
La Trobe University, Technical Report, 2001.

\bibitem{seccomp}
J. Edge, ``A seccomp overview,'' \textit{LWN.net}, 2015.

\bibitem{gvisor}
Google, ``gVisor: Application Kernel for Containers,''
\textit{USENIX ATC}, 2018.

\bibitem{code-sandbox}
M. Christodorescu et al., ``Towards Secure AI-Generated Code,'' \textit{arXiv
  preprint}, 2024.

\bibitem{codex}
M. Chen et al., ``Evaluating Large Language Models Trained on Code,''
\textit{arXiv:2107.03374}, 2021.

\bibitem{code-llm-survey}
Z. Zheng et al., ``A Survey of Large Language Models for Code,''
\textit{arXiv:2311.08947}, 2023.

\bibitem{llm-compression}
T. Dettmers et al., ``QLoRA: Efficient Finetuning of Quantized LLMs,''
\textit{NeurIPS}, 2023.

\bibitem{wasm-generation}
B. Bichsel et al., ``CodeGen2: Lessons for Enterprise-Grade Code Generation,''
\textit{arXiv:2305.02309}, 2023.

\end{thebibliography}

\end{document}
